\section{Related Work}
\tempo{} is rooted in the reactive control lineage: Esterel, FairThreads, and
ReactiveML. It inherits Boussinot's delayed-absence discipline
\cite{esterel1992,fairthreads2003,reactiveml2005}. Relative to Esterel, \tempo{}
does not impose a static compilation discipline; instead it targets dynamic
behavior and library-level composition. Relative to FairThreads, \tempo{} provides
a typed OCaml API with higher-order combinators and valued signals.

Other synchronous languages in the same family include Lustre and Signal, which
take dataflow-oriented approaches to synchronous programming, and the broader
survey literature that situates these languages in a shared semantic landscape
\cite{lustre1991,signal1987,benveniste2003}. \tempo{} aligns with this lineage but
focuses on control primitives realized at library level rather than a dedicated
compiler.

Industrial synchronous formalisms have also shaped this landscape. Statecharts
provide a widely adopted visual model for reactive control
\cite{harel1987statecharts}, and SCADE offers a synchronous design and
validation environment for safety-critical embedded control software
\cite{berry2007scade}. These systems emphasize predictable control behavior and
tooling, whereas \tempo{} targets a lightweight, library-level embedding.

Closely related academic languages include Lucid Synchrone and Quartz, which
extend synchronous ideas with functional or model-based design perspectives
\cite{caspihamonpouzet2008lucid,schneider2017quartz}. These reinforce the
continuity between \tempo{} and earlier synchronous language research.

ReactiveML is the closest antecedent: it integrates reactive constructs into a
language extension with dedicated syntax and compilation support. \tempo{}, by
contrast, is a pure library that relies on OCaml 5.3 algebraic effects to
implement the same control primitives without modifying the compiler or the
runtime. This makes \tempo{} easier to integrate into existing OCaml codebases and
enables experimentation with scheduling policies in user space, at the cost of
foregoing some static analyses and language-level guarantees provided by
ReactiveML.

In contrast with FRP \cite{frp1997}, \tempo{} emphasizes control and instant-based
determinism rather than continuous dataflow. The library focuses on explicit
signals, guards, and preemption, aligning with synchronous reactive semantics
instead of stream processing abstractions.

Stream-processing languages such as Lucid and StreamIt
\cite{ashcroft1977lucid,thies2002streamit} emphasize dataflow pipelines rather
than control-centric instants. They provide a useful contrast to \tempo{}: the
dataflow style is oriented toward throughput and steady-state streaming, while
\tempo{} targets explicit control reactions at discrete instants.

Recent FRP systems have introduced stronger operational guarantees. The
modal FRP line of work uses modal types to enforce causality and prevent space
leaks, including Simply RaTT and its extensions
\cite{bahr2019ratt,bahr2022modalfrp,bahr2023async}. These systems target
dataflow FRP rather than control-oriented primitives, but they offer insights
into static enforcement mechanisms that could complement \tempo{}'s current
runtime checks.
Recent applications of modal FRP include GUI programming with modal types,
which demonstrates the practicality of these systems in reactive user
interfaces \cite{bahr2025gui}.

Industrial reactive-streams frameworks (e.g., RxJava, Akka Streams) emphasize
asynchronous stream processing and backpressure rather than synchronous
instants. While their semantics differ from \tempo{}'s model, they represent a
separate lineage of reactive programming practice and are surveyed in the
reactive-streams literature \cite{cheng2017reactivestreams,davis2019reactivestreams}.

\tempo{} also builds on the recent availability of algebraic effects and handlers
in OCaml \cite{multicoreocaml2020}. At a semantic level, effects provide a
structured interface to delimited continuations
\cite{danvyfilinski1990,filinski1994,plotkinpretnar2009}, which has long been
recognized as a natural foundation for user-level schedulers and cooperative
threads. Systems such as Eff demonstrated the practicality of effect handlers
as a runtime control mechanism \cite{eff2013}, and \tempo{} leverages the same
mechanism for synchronous control, rather than for general-purpose effects.

Compared with effect-based concurrency libraries, \tempo{}'s distinctive
contribution is its \emph{instant} discipline: continuations are not resumed
arbitrarily, but at well-defined boundaries that expose synchronous semantics.
This positions \tempo{} between language-level synchronous systems (Esterel,
ReactiveML) and general-purpose effect-based schedulers, offering a library
path to control-oriented reactive programming.
